\chapter{Use Case Modeling}
\section{Use Case Diagram}
The workbook contains 11 actors. The principal human roles are Driver, Gate Attendant, Parking Administrator, and Parking Manager. Entry/exit cameras, slot sensors, fire alarms, the payment gateway, charging stations, and the central server are supporting actors. The diagrams are grouped into three figures for readability: Figures~\ref{fig:uc-driver}, \ref{fig:uc-exit}, and \ref{fig:uc-admin}. Every one of the 22 use cases appears in at least one of them (UC-12 appears in two), and the index in Table~\ref{tab:uc-index} lists all of them.
\tikzset{ucs/.style={uc,text width=21mm},lbl/.style={font=\tiny,fill=white,inner sep=1pt}}
\widefigure{\begin{tikzpicture}[x=1cm,y=1cm]
\node[ext] (driver) at (0.5,-6) {Driver};
\node[ucs] (u01) at (5.4,0) {UC-01\\Vehicle entry};
\node[ucs] (u18) at (5.4,-2) {UC-18\\Pass / VIP entry};
\node[ucs] (u20) at (5.4,-4) {UC-20\\EV charging};
\node[ucs] (u03) at (5.4,-6) {UC-03\\Reserve slot};
\node[ucs] (u13) at (5.4,-8) {UC-13\\Cancel reservation};
\node[ucs] (u04) at (5.4,-10) {UC-04\\Locate vehicle};
\node[ucs] (u15) at (5.4,-12) {UC-15\\View availability};
\node[ucs] (u07) at (10.6,0) {UC-07\\Read plate};
\node[ucs] (u08) at (10.6,-2) {UC-08\\Nearest slot};
\node[ext] (cam) at (15.8,0) {Entry Camera};
\node[ext] (charger) at (15.8,-4) {Charging Station};
\foreach \u in {u01,u18,u20,u03,u13,u04,u15} \draw (driver)--(\u);
\draw (cam)--(u07);\draw (charger)--(u20);
\draw[arr,dashed] (u01)--node[lbl]{include}(u07);
\draw[arr,dashed] (u01)--node[lbl,pos=.3]{include}(u08);
\draw[arr,dashed] (u18)--node[lbl,pos=.75]{include}(u07);
\draw[arr,dashed] (u20)--node[lbl,pos=.55]{include}(u08);
\draw[arr,dashed] (u18)--node[lbl]{extend}(u01);
\begin{scope}[on background layer]\node[draw=muted,dashed,fit=(u01)(u15)(u07)(u08),inner sep=5mm,label={[font=\small]above:SPS -- Entry, reservation, and information}] {};\end{scope}
\end{tikzpicture}}{Driver use cases around entry, reservation, and information (UC-01, 03, 04, 07, 08, 13, 15, 18, 20)}{fig:uc-driver}

\widefigure{\begin{tikzpicture}[x=1cm,y=1cm]
\node[ext] (driver) at (0.5,0) {Driver};
\node[ext] (pay) at (0.5,-3) {Payment Gateway};
\node[ucs] (u02) at (5.4,0) {UC-02\\Exit and pay};
\node[ucs] (u12) at (5.4,-4) {UC-12\\Log in};
\node[ucs] (u21) at (10.6,2) {UC-21\\Verify exit plate};
\node[ucs] (u11) at (10.6,0) {UC-11\\Dynamic fee};
\node[ucs] (u06) at (10.6,-2) {UC-06\\Lost ticket};
\node[ucs] (u17) at (10.6,-4) {UC-17\\Override fee};
\node[ext] (xcam) at (15.8,2) {Exit Camera};
\node[ext] (att) at (15.8,-2) {Gate Attendant};
\node[ext] (mgr) at (15.8,-4) {Parking Manager};
\draw (driver)--(u02);\draw (pay)--(u02);\draw (xcam)--(u21);
\draw (att)--(u06);\draw (att)--(u17);\draw (mgr)--(u17);\draw (mgr)--(u06);
\draw[arr,dashed] (u02)--node[lbl]{include}(u21);
\draw[arr,dashed] (u11)--node[lbl]{extend}(u02);
\draw[arr,dashed] (u06)--node[lbl,pos=.4]{extend}(u02);
\draw[arr,dashed] (u17)--node[lbl,pos=.3]{extend}(u02);
\draw[arr,dashed] (u17)--node[lbl]{include}(u12);
\begin{scope}[on background layer]\node[draw=muted,dashed,fit=(u02)(u12)(u21)(u17),inner sep=5mm,label={[font=\small]above:SPS -- Exit and payment}] {};\end{scope}
\end{tikzpicture}}{Exit, payment, and exceptional fee handling (UC-02, 06, 11, 12, 17, 21)}{fig:uc-exit}

\widefigure{\begin{tikzpicture}[x=1cm,y=1cm]
\node[ext] (admin) at (0.5,-2.25) {Administrator};
\node[ext] (mgr) at (0.5,-6) {Parking Manager};
\node[ext] (att) at (0.5,-8.25) {Gate Attendant};
\node[ucs] (u10) at (5.4,0) {UC-10\\Evacuation};
\node[ucs] (u05) at (5.4,-1.5) {UC-05\\Lock zone};
\node[ucs] (u14) at (5.4,-3) {UC-14\\Configure lot};
\node[ucs] (u19) at (5.4,-4.5) {UC-19\\Register pass};
\node[ucs] (u16) at (5.4,-6) {UC-16\\Report};
\node[ucs] (u22) at (5.4,-7.5) {UC-22\\Work offline};
\node[ucs] (u09) at (5.4,-9) {UC-09\\Monitor slots};
\node[ucs] (u12) at (10.6,-3.75) {UC-12\\Log in};
\node[ext] (fire) at (15.8,0) {Fire Alarm};
\node[ext] (central) at (15.8,-6) {Central Server};
\node[ext] (sens) at (15.8,-9) {Slot Sensor};
\foreach \u in {u10,u05,u14,u19} \draw (admin)--(\u);
\draw (mgr)--(u16);\draw (att)--(u22);\draw (att)--(u09);
\draw (fire)--(u10);\draw (sens)--(u09);\draw (central)--(u19);\draw (central)--(u22);
\foreach \u in {u05,u14,u19,u16} \draw[arr,dashed] (\u)--node[lbl,pos=.65]{include}(u12);
\begin{scope}[on background layer]\node[draw=muted,dashed,fit=(u10)(u09)(u12),inner sep=5mm,label={[font=\small]above:SPS -- Administration, safety, and operations}] {};\end{scope}
\end{tikzpicture}}{Administration, safety, and offline operation (UC-05, 09, 10, 12, 14, 16, 19, 22)}{fig:uc-admin}

An \emph{include} relationship represents reused behavior and is drawn from the base use case to the included one: UC-01 includes UC-07 and UC-08; UC-02 includes UC-21; UC-18 includes UC-07; UC-20 includes UC-08; UC-05, UC-14, UC-16, UC-17, and UC-19 include UC-12. An \emph{extend} relationship is drawn from the extending use case to the base one: UC-18 (monthly/VIP pass) extends UC-01, and UC-11 (dynamic fee), UC-17 (fee override), and UC-06 (lost ticket) extend UC-02. A fire-alarm signal can start UC-10 without waiting for an operator login.
\par\begingroup\footnotesize\setlength{\tabcolsep}{4pt}
\begin{longtable}{>{\raggedright\arraybackslash}p{\dimexpr 0.11000\linewidth-2\tabcolsep\relax}>{\raggedright\arraybackslash}p{\dimexpr 0.50000\linewidth-2\tabcolsep\relax}>{\raggedright\arraybackslash}p{\dimexpr 0.39000\linewidth-2\tabcolsep\relax}}
\caption{Complete index of 22 use cases}\label{tab:uc-index}\\
\toprule
\textbf{ID} & \textbf{Use case / actor} & \textbf{Related requirements} \\ \midrule
\endfirsthead
\multicolumn{3}{l}{\small\itshape Complete index of 22 use cases (continued)}\\
\toprule
\textbf{ID} & \textbf{Use case / actor} & \textbf{Related requirements} \\ \midrule
\endhead
\midrule
\multicolumn{3}{r}{\scriptsize Continued on the next page}\\
\endfoot
\bottomrule\endlastfoot
UC-01 & Vehicle entry / Driver & REQ-01; REQ-02; REQ-31 \\
UC-02 & Exit and payment / Driver & REQ-03; REQ-04; REQ-26; REQ-31 \\
UC-03 & Reserve a slot / Driver & REQ-05 \\
UC-04 & Locate a vehicle / Driver & REQ-06 \\
UC-05 & Lock or unlock a zone / Administrator & REQ-08; REQ-09; REQ-18 \\
UC-06 & Handle a lost ticket / Gate Attendant & REQ-07 \\
UC-07 & Recognize an entry plate / Driver / Entry Camera & REQ-11; REQ-01 \\
UC-08 & Assign the nearest slot / Driver & REQ-12 \\
UC-09 & Detect parking discrepancies / Gate Attendant / Slot Sensor & REQ-13; REQ-14 \\
UC-10 & Emergency evacuation / Administrator / Fire Alarm & REQ-15; REQ-09 \\
UC-11 & Pay a dynamic fee / Driver & REQ-16; REQ-17; REQ-27; REQ-28 \\
UC-12 & Log in / Attendant / Manager / Administrator & REQ-18; REQ-09 \\
UC-13 & Cancel a reservation / Driver & REQ-19; REQ-05 \\
UC-14 & Configure lot and rates / Administrator & REQ-20; REQ-09 \\
UC-15 & View availability / Driver & REQ-02 \\
UC-16 & Revenue and audit report / Manager & REQ-21; REQ-10; REQ-09 \\
UC-17 & Override a fee / Attendant / Manager & REQ-22; REQ-03; REQ-09 \\
UC-18 & Use a monthly/VIP pass / Driver & REQ-23; REQ-12 \\
UC-19 & Register/remove a pass / Administrator & REQ-23; REQ-09; REQ-29 \\
UC-20 & Charge an EV / Driver / Charging Station & REQ-24; REQ-25 \\
UC-21 & Verify an exit plate / Driver / Exit Camera & REQ-26; REQ-11 \\
UC-22 & Work offline / Gate Attendant / Central Server & REQ-29; REQ-30; REQ-15 \\
\end{longtable}\endgroup


\section{Detailed Scenarios}
Each specification records the actor, preconditions, main flow, exceptions, postconditions, and linked requirements. The scenarios emphasize business decisions. Original fields, frequencies, triggers, and complete source text are available in \texttt{data/dataset.json}. Corresponding execution results are stored in \texttt{evidence/use\_case\_runs.txt}.
\subsection{UC-01 -- Vehicle entry}
\noindent\textbf{Actor:} Driver\quad\textbf{Traceability:} REQ-01; REQ-02; REQ-31.
\par\textbf{Preconditions:} A matching slot is available; the lot is not evacuating.
\par\textbf{Main flow:} Read and confirm the plate; select a reserved or nearest matching slot; issue a ticket and guidance; open and close the entry gate.
\par\textbf{Alternative flows and exceptions:} Reject full lots, evacuation, and duplicate plates. Low confidence requires manual input.
\par\textbf{Postconditions:} Ticket Active; slot Occupied; entry logged.
\subsection{UC-02 -- Exit and payment}
\noindent\textbf{Actor:} Driver\quad\textbf{Traceability:} REQ-03; REQ-04; REQ-26; REQ-31.
\par\textbf{Preconditions:} An active ticket exists.
\par\textbf{Main flow:} Scan the ticket; verify the plate; calculate and display the fee; take payment; close the ticket, release the slot, and operate the exit gate.
\par\textbf{Alternative flows and exceptions:} Missing ticket invokes lost-ticket handling. Failed payment retains the ticket; permit retry or another method. A plate mismatch needs confirmation.
\par\textbf{Postconditions:} Ticket Closed; slot Free or Blocked if locked; revenue recorded.
\subsection{UC-03 -- Reserve a slot}
\noindent\textbf{Actor:} Driver\quad\textbf{Traceability:} REQ-05.
\par\textbf{Preconditions:} A matching slot exists; the central server is reachable.
\par\textbf{Main flow:} Enter type, plate, and arrival time; find a slot; mark Reserved and confirm the hold.
\par\textbf{Alternative flows and exceptions:} Reject if full. Release a no-show after the 15-minute hold when expiry is checked.
\par\textbf{Postconditions:} Reservation recorded; slot unavailable to other vehicles.
\subsection{UC-04 -- Locate a vehicle}
\noindent\textbf{Actor:} Driver\quad\textbf{Traceability:} REQ-06.
\par\textbf{Preconditions:} The vehicle has an active ticket in the lot.
\par\textbf{Main flow:} Enter the plate; look up the ticket; return floor, zone, slot, and distance.
\par\textbf{Alternative flows and exceptions:} If absent from active tickets, display vehicle not found.
\par\textbf{Postconditions:} Location displayed without modifying data.
\subsection{UC-05 -- Lock or unlock a zone}
\noindent\textbf{Actor:} Administrator\quad\textbf{Traceability:} REQ-08; REQ-09; REQ-18.
\par\textbf{Preconditions:} An administrator account exists; the lot is in a suitable mode.
\par\textbf{Main flow:} Log in; select a floor/zone; lock or unlock; update Free/Blocked states and log the action.
\par\textbf{Alternative flows and exceptions:} Deny insufficient privileges. Occupied slots remain until exit, then become Blocked if locked. Faulty slots remain blocked on unlock.
\par\textbf{Postconditions:} The locked area cannot take new visitors; existing reservations follow the stated policy.
\subsection{UC-06 -- Handle a lost ticket}
\noindent\textbf{Actor:} Gate Attendant\quad\textbf{Traceability:} REQ-07.
\par\textbf{Preconditions:} A driver reports a lost ticket and provides a plate.
\par\textbf{Main flow:} Verify ownership; find the ticket; calculate the charge; close the ticket, release the slot, and log the incident.
\par\textbf{Alternative flows and exceptions:} No record: escalate to the manager. No ownership proof: refuse exit and record the denial.
\par\textbf{Postconditions:} Ticket Closed; charge and incident recorded; REQ-07 criteria still require completion.
\subsection{UC-07 -- Recognize an entry plate}
\noindent\textbf{Actor:} Driver / Entry Camera\quad\textbf{Traceability:} REQ-11; REQ-01.
\par\textbf{Preconditions:} The vehicle is stopped at entry; recognition or manual input is available.
\par\textbf{Main flow:} The simulated camera returns plate and confidence; accept at least 80\%; pass the confirmed plate to entry.
\par\textbf{Alternative flows and exceptions:} Below 80\% or camera offline: request manual input and log the reason.
\par\textbf{Postconditions:} A confirmed plate is available before ticket issuance.
\subsection{UC-08 -- Assign the nearest slot}
\noindent\textbf{Actor:} Driver\quad\textbf{Traceability:} REQ-12.
\par\textbf{Preconditions:} A suitable slot is available.
\par\textbf{Main flow:} Filter by type, state, and VIP eligibility; select minimum distance; break ties by tree order; display the route.
\par\textbf{Alternative flows and exceptions:} Skip locked/full nearest zones; return no slot when suitable capacity is exhausted.
\par\textbf{Postconditions:} A valid slot is selected for ticket issuance.
\subsection{UC-09 -- Detect parking discrepancies}
\noindent\textbf{Actor:} Gate Attendant / Slot Sensor\quad\textbf{Traceability:} REQ-13; REQ-14.
\par\textbf{Preconditions:} Sensor readings and ticket data are available.
\par\textbf{Main flow:} Read each slot; reconcile occupancy with tickets; alert and log; relocate a ticket if staff confirm the vehicle is in another slot.
\par\textbf{Alternative flows and exceptions:} No heartbeat makes the slot Blocked. Stays over 24 hours trigger an overstay alert. An unticketed vehicle triggers an unauthorized-vehicle alert.
\par\textbf{Postconditions:} Discrepancies recorded; faulty slots cannot receive new vehicles.
\subsection{UC-10 -- Emergency evacuation}
\noindent\textbf{Actor:} Administrator / Fire Alarm\quad\textbf{Traceability:} REQ-15; REQ-09.
\par\textbf{Preconditions:} An administrator command or fire-alarm signal is received.
\par\textbf{Main flow:} Open all barriers; stop tickets; cancel reservations and lock the lot; waive exit fees; let an administrator end and restore operation.
\par\textbf{Alternative flows and exceptions:} Deny attendant requests to trigger through administration or to end evacuation. A fire alarm bypasses login.
\par\textbf{Postconditions:} Barriers stay open during evacuation; completion restores Normal mode.
\subsection{UC-11 -- Pay a dynamic fee}
\noindent\textbf{Actor:} Driver\quad\textbf{Traceability:} REQ-16; REQ-17; REQ-27; REQ-28.
\par\textbf{Preconditions:} An active ticket exists.
\par\textbf{Main flow:} Split the stay into started hours; apply peak multipliers; cap each 24-hour block; apply the ticket surcharge; display and collect the fee.
\par\textbf{Alternative flows and exceptions:} Grace-period or monthly/VIP stays waive parking. Without peak hours or surcharge, apply base pricing.
\par\textbf{Postconditions:} The exit workflow records a fee consistent with the rules.
\subsection{UC-12 -- Log in}
\noindent\textbf{Actor:} Attendant / Manager / Administrator\quad\textbf{Traceability:} REQ-18; REQ-09.
\par\textbf{Preconditions:} The account exists and is not locked.
\par\textbf{Main flow:} Enter credentials; authenticate; return the role; establish the session and record the event.
\par\textbf{Alternative flows and exceptions:} Wrong credentials increment the count; three consecutive failures lock the account; locked accounts cannot log in.
\par\textbf{Postconditions:} A role is returned or access is denied.
\subsection{UC-13 -- Cancel a reservation}
\noindent\textbf{Actor:} Driver\quad\textbf{Traceability:} REQ-19; REQ-05.
\par\textbf{Preconditions:} A reservation exists for the plate.
\par\textbf{Main flow:} Find the reservation; confirm cancellation; clear the hold; set Free or Blocked if locked; log the action.
\par\textbf{Alternative flows and exceptions:} Show distinct messages for missing and expired reservations.
\par\textbf{Postconditions:} The slot can be reallocated when eligible.
\subsection{UC-14 -- Configure lot and rates}
\noindent\textbf{Actor:} Administrator\quad\textbf{Traceability:} REQ-20; REQ-09.
\par\textbf{Preconditions:} An administrator account exists.
\par\textbf{Main flow:} Log in; select rate change or new zone; validate; apply and log.
\par\textbf{Alternative flows and exceptions:} Reject nonpositive rates, invalid types, duplicates, or insufficient privileges. Existing tickets retain their rates.
\par\textbf{Postconditions:} New tickets use the new rate; new zones are usable in the current session.
\subsection{UC-15 -- View availability}
\noindent\textbf{Actor:} Driver\quad\textbf{Traceability:} REQ-02.
\par\textbf{Preconditions:} The lot hierarchy is available.
\par\textbf{Main flow:} Traverse lot, floors, and zones; count Free leaves; display free and total counts.
\par\textbf{Alternative flows and exceptions:} Locked or evacuating branches cannot accept new allocations; compare counts with flags and states.
\par\textbf{Postconditions:} Query does not modify data.
\subsection{UC-16 -- Revenue and audit report}
\noindent\textbf{Actor:} Manager\quad\textbf{Traceability:} REQ-21; REQ-10; REQ-09.
\par\textbf{Preconditions:} Logged in as manager or administrator.
\par\textbf{Main flow:} Select a period; read payments and lost-ticket charges; group by zone/floor; return totals and audit events.
\par\textbf{Alternative flows and exceptions:} No transactions gives zero revenue. Deny attendants access.
\par\textbf{Postconditions:} Report displayed without modifying source records.
\subsection{UC-17 -- Override a fee}
\noindent\textbf{Actor:} Attendant / Manager\quad\textbf{Traceability:} REQ-22; REQ-03; REQ-09.
\par\textbf{Preconditions:} An active ticket and an authorized requesting role exist.
\par\textbf{Main flow:} Enter plate, amount, and reason; create Pending request; manager/admin approves; set the replacement amount and log the decision.
\par\textbf{Alternative flows and exceptions:} Reject negative amounts, missing reasons, or attendant approval. Rejection retains the calculated fee.
\par\textbf{Postconditions:} The override is effective only after approval.
\subsection{UC-18 -- Use a monthly/VIP pass}
\noindent\textbf{Actor:} Driver\quad\textbf{Traceability:} REQ-23; REQ-12.
\par\textbf{Preconditions:} A plate-based pass is valid at entry.
\par\textbf{Main flow:} Recognize the plate; find the pass; prioritize VIP slots for VIP holders; record pass type on the ticket; waive parking at exit.
\par\textbf{Alternative flows and exceptions:} Expired passes use ordinary pricing. Monthly passes cannot use VIP slots. EV energy remains chargeable.
\par\textbf{Postconditions:} Exit closes the ticket; the accepted pass makes parking free.
\subsection{UC-19 -- Register/remove a pass}
\noindent\textbf{Actor:} Administrator\quad\textbf{Traceability:} REQ-23; REQ-09; REQ-29.
\par\textbf{Preconditions:} Administrator privileges and central-server connectivity exist.
\par\textbf{Main flow:} Log in; enter plate, monthly/vip type, and expiry date; validate, save, and log; support pass removal.
\par\textbf{Alternative flows and exceptions:} Reject unknown types, past dates, disconnection, or insufficient privileges.
\par\textbf{Postconditions:} Pass registered or removed in the simulation data.
\subsection{UC-20 -- Charge an EV}
\noindent\textbf{Actor:} Driver / Charging Station\quad\textbf{Traceability:} REQ-24; REQ-25.
\par\textbf{Preconditions:} An EV slot and valid EV ticket are available.
\par\textbf{Main flow:} Allocate the nearest EV slot; receive and accumulate kWh; charge parking plus energy on exit.
\par\textbf{Alternative flows and exceptions:} Reject when EV slots are full; do not allocate EV slots to cars; reject nonpositive kWh.
\par\textbf{Postconditions:} Energy is included in the EV ticket total.
\subsection{UC-21 -- Verify an exit plate}
\noindent\textbf{Actor:} Driver / Exit Camera\quad\textbf{Traceability:} REQ-26; REQ-11.
\par\textbf{Preconditions:} A ticket and camera plate or explicit unreadable result are available.
\par\textbf{Main flow:} Read the exit plate; find the ticket; compare; proceed to payment on a match.
\par\textbf{Alternative flows and exceptions:} Mismatch/unreadable: hold the barrier, log, and call staff; authorized confirmation may release the vehicle.
\par\textbf{Postconditions:} Release or hold is recorded; omitted input remains an implementation limitation.
\subsection{UC-22 -- Work offline}
\noindent\textbf{Actor:} Gate Attendant / Central Server\quad\textbf{Traceability:} REQ-29; REQ-30; REQ-15.
\par\textbf{Preconditions:} Local data remains usable.
\par\textbf{Main flow:} Detect disconnection; switch to local mode; allow entry/exit with simulated cash/card; mark logs; process marks and record their count on reconnection.
\par\textbf{Alternative flows and exceptions:} Refuse online payment, reservations, and pass updates. Evacuation remains available offline.
\par\textbf{Postconditions:} Local transactions remain in audit; actual network synchronization is not implemented.


\section{Chapter Conclusion}
Use cases translate requirements into executable interactions, especially for capacity shortages, camera faults, payment failure, plate mismatches, and connection loss. Access conditions and evacuation priority appear alongside normal flows so that validation extends beyond successful cases.

The 22 use cases are drawn in three diagrams and specified with preconditions, a main flow, alternative flows tagged by step, and postconditions. The include and extend relations are written in the flow text, so each relation can be read and tested in the same place. Every functional requirement has at least one use-case link (Table~\ref{tab:coverage}), and the use cases are run on the simulation, which turns the descriptions into repeatable evidence. What the use cases do not show is hardware behavior: cameras, sensors, and payment are simulated, so a successful scenario demonstrates the business logic only.
